\section{Experiments}
\label{sec:experiments}

\subsection{Baseline Models}
\label{sec:baselines}

We evaluate two open-weight models in a zero-shot repair setting:
\begin{itemize}
    \item \textbf{DeepSeek-R1 1.5B}: A small reasoning model, serving as a performance floor.
    \item \textbf{Qwen2.5-Coder 7B}: A code-specialized model, serving as the main baseline.
\end{itemize}

Both models are run locally via Ollama. The repair prompt provides the broken configuration and its validation errors, asking the model to return only the fixed configuration.

\subsection{Results}
\label{sec:results}

% TODO: Insert results tables (auto-generated from baseline runs)

% \begin{table}[h]
% \centering
% \caption{Overall pass@$k$ results.}
% \label{tab:overall_results}
% \begin{tabular}{lcc}
% \toprule
% Model & pass@1 & pass@3 \\
% \midrule
% DeepSeek-R1 1.5B & -- & -- \\
% Qwen2.5-Coder 7B & -- & -- \\
% \bottomrule
% \end{tabular}
% \end{table}

\subsection{Analysis}
\label{sec:analysis}

% TODO: Add analysis by category, difficulty, and format
% - Syntactic faults are easiest to repair
% - Semantic and cross-resource faults are most challenging
% - CF vs TF comparison
% - Difficulty level correlation with repair success
